\section{Discussion}
\label{sec:discussion}

\textbf{Workload-dependent benefits.}
Our experiments show that the magnitude of \mysys's benefit depends strongly on
workload characteristics. Recommendation and decode workloads benefit most
because Attention and MatMul expose profitable invariant schedules, yielding
about 10\% average inference speedup. Schedule-fixed expert-written libraries
also benefit from removing unreachable template instances. This per-model
pruning is most attractive when a deployment serves fewer models, because each
pruned binary can discard a larger fraction of the universal library. By
contrast, HSTU and prefill remain compute-heavy, so specializing their core
Attention and MatMul kernels removes little of the total work and improves
end-to-end latency by only about 3\%. Some kernels in the current open-source
baselines also cannot be further specialized by \mysys, further limiting
end-to-end gains.

\textbf{Platform applicability.}
Our schedule-generic kernel-specialization implementation currently targets
Ascend. In the open-source libraries we surveyed for other accelerators,
expert-written kernels predominantly follow the schedule-fixed design, leaving
no comparable schedule-generic baseline to evaluate. Ascend exposes many
programmable hardware levels; achieving high performance therefore relies on
numerous runtime schedule parameters, which motivates schedule-generic
libraries and makes automatic specialization particularly relevant.

\textbf{Optimization scope and limitations.}
\mysys applies to existing expert-written operator libraries: it specializes
schedule-generic kernels or removes unreachable instances from schedule-fixed
libraries. This differs from AI compilers, which generate implementations and
search schedules; their code-generation pipelines do not optimize an existing
C++ expert library in place. 

In addition, \mysys analyzes host code only to recover constants
and launch-site liveness, then optimizes reachable device kernels; it does not
rewrite host code. Host control flow is complex, CPUs already execute it
efficiently, and profitable host specialization would require a more selective
policy than device-kernel constant propagation. NPU Graph replay suppresses
repeated launch overhead, while identical input shapes across Transformer
layers allow host-computed schedule parameters to be reused. Host schedule
computation is therefore not a bottleneck in our evaluated setting, making
device-kernel optimization the higher-value target.
